\section{Related Work}

\paragraph{Normal-only video anomaly detection.}
The established VAD setting fits a model of ordinary activity and evaluates deviations at test time. Early crowded-scene work combined dynamic texture models and discriminative saliency on UCSD data \citep{li2014anomaly}; Avenue was introduced with an efficient sparse-combination learning method \citep{lu2013abnormal}. Deep methods later used future-frame prediction to make abnormal motion or appearance difficult to predict \citep{liu2018future}. Memory-based approaches restrict reconstruction to learned normal prototypes, either by adding a memory to an autoencoder \citep{gong2019memorizing} or by coupling memory items to compactness and separateness losses \citep{park2020memory}. Our study does not compare end-to-end architectures. It freezes one general visual representation and asks whether low-capacity normality scorers fail because that representation contains dataset/camera nuisance information.

\paragraph{Scene dependence and evaluation.}
Anomalous behavior is contextual: an action may be ordinary in one location and exceptional in another. The Street Scene benchmark emphasized that single-scene and multi-scene formulations define different tasks and proposed criteria intended to better reflect practical event detection \citep{ramachandra2020street}. This observation creates a tension for scene invariance. Removing camera identity may improve transfer, but aggressive removal can also discard the location information that makes an event anomalous. SRN represents this tension explicitly through a constrained residual branch and a small context branch. Our evaluation therefore measures both dataset/camera decodability and anomaly ranking instead of interpreting lower scene accuracy alone as success.

\paragraph{Cross-domain VAD.}
Cross-domain work considers a detector trained on one collection and evaluated on another. For example, zxVAD studies cross-domain inference without target-domain adaptation and uses pseudo-anomalies plus a Normalcy Classifier module with relative-normalcy learning \citep{aich2023crossdomain}. That setting and ours are not directly comparable: zxVAD learns a video prediction model with auxiliary data, whereas we test normal-only scorers on a shared frozen cache and prohibit anomaly labels during selection. The methodological point we retain is that same-dataset and cross-dataset results answer different questions. We add a threshold distinction inside the latter: AUROC describes ordering after observing the full test curve, while a source-fixed threshold tests whether the numerical score scale transfers.

Multi-domain VAD has also been formalized through domain-generalization methods and broader benchmarks. Cho et al. formalize MDVAD, including leave-one-out training on multiple source domains and evaluation on an unseen target domain \citep{cho2024multidomain}, while MSAD introduces a 14-scenario multi-scenario benchmark \citep{zhu2024advancing}. Our two-dataset study cannot reproduce that evidential scope. Its role is a controlled falsification pilot for one residualization mechanism, not a multi-domain benchmark result.

\paragraph{Invariant and frozen representations.}
Domain-adversarial training uses a gradient-reversal layer to discourage domain prediction while optimizing a task representation \citep{ganin2016domain}. We include this construction as a generic residualization control and also apply it inside SRN. DINOv2 supplies the common frozen input because its self-supervised training targets broadly useful visual features \citep{oquab2024dinov2}. Freezing the backbone removes a major capacity confound: every scorer and residual variant receives exactly the same 384-dimensional frame feature. It also limits the conclusion. We study what a small residual module can do to one frozen image representation; we do not infer that temporal, object-centric, or end-to-end domain-generalization methods must behave similarly.

The closest evaluation precedent is Rashidi's cross-dataset audit of off-the-shelf frozen representations, including DINOv2, kNN, a Mahalanobis control, and false alarms per hour \citep{rashidi2026benchmark}. We do not claim novelty for auditing frozen-feature transfer. We instead add a predeclared source-threshold versus target-normal-calibration split and use a capacity-matched residualization matrix plus an independent identity probe to test whether a specific proposed mechanism repairs the observed shift.

\paragraph{Calibration and operating points.}
VAD papers commonly report frame AUROC, yet surveillance workloads care about low false-positive operation and alarm burden. Street Scene's evaluation critique motivates reporting explicit false-positive burden beyond frame AUROC; we therefore include false-alarm events per hour where frame rate is known \citep{ramachandra2020street}. We further separate two admissible thresholds. The strict threshold is estimated only from held-out source-normal scores. The calibrated threshold uses a predeclared set of target-normal videos and never uses target anomalies. Neither is the oracle threshold implicit in TPR at a selected point of the final test ROC curve. This separation is central to the empirical finding rather than a post-hoc metric choice.
